\begin{textsample}{POS Dim 5 – ro – Score 10.00 – ro/ro\_ef\_52.txt}  \label{ex:f5_pos_115}
{\EmojiFont ➢} Relacionar o texto com suas condições de produção, seu contexto sócio-histórico de circulação e com os projetos de dizer : leitor e leitura previstos, objetivos, pontos de vista e perspectivas em jogo, papel social do autor, época, gênero do discurso e esfera/campo em questão etc. {\EmojiFont ➢} Analisar a circulação dos gêneros do discurso nos diferentes campos de atividade, seus usos e funções relacionados com as atividades típicas do campo, seus diferentes agentes, os interesses em jogo e as práticas de linguagem em circulação e as relações de determinação desses elementos sobre a construção composicional, as marcas linguísticas ligadas ao estilo e o conteúdo temático dos gêneros. {\EmojiFont ➢} Refletir sobre as transformações ocorridas nos campos de atividades em função do desenvolvimento das tecnologias de comunicação e informação, do uso do \textbf{hipertexto} e da hipermídia e do surgimento da Web 2.0 : novos gêneros do discurso e novas práticas de linguagem próprias da cultura digital, transmutação ou reelaboração dos gêneros em função das transformações pelas quais passam o texto ( de formatação e em função da convergência de mídias e do funcionamento hipertextual ), novas formas de interação e de compartilhamento de textos/conteúdos/informações, reconfiguração do papel de leitor, que passa a ser também produtor, dentre outros, como forma de ampliar as possibilidades de participação na cultura digital e contemplar os novos e os \textbf{multiletramentos}. {\EmojiFont ➢} Fazer apreciações e valorações estéticas, éticas, políticas e ideológicas, dentre outras, envolvidas na leitura crítica de textos verbais e de outras produções culturais.
Reconstrução e reflexão sobre as condições de produção e recepção dos textos pertencentes a diferentes gêneros e que circulam nas diferentes mídias e esferas/campos de atividade humana {\EmojiFont ➢} Analisar as diferentes formas de manifestação da compreensão ativa ( réplica ativa ) dos textos que circulam nas redes sociais, blogs/microblog, sites e afins e os gêneros que conformam essas práticas de linguagem, como : comentário, carta de leitor, post em rede social, gif, meme, fanfic, vlogs variados, political remix, charge digital, \textbf{paródias} de diferentes tipos, vídeos-minuto, e-zine, fanzine, fanvídeo, vidding, gameplay, walkthroug, detonado, machinima, trailer honesto, playlists \textbf{comentadas} de diferentes tipos etc., de forma a ampliar a compreensão de textos que pertencem a esses gêneros e a possibilitar uma participação mais qualificada do ponto de vista ético, estético e político nas práticas de linguagem da cultura digital.
Dialogia e relação entre textos {\EmojiFont ➢} Identificar e refletir sobre as diferentes perspectivas ou vozes presentes nos textos e sobre os efeitos de sentido do uso do discurso direto, indireto, indireto livre, citações etc. {\EmojiFont ➢} Estabelecer relações de intertextualidade e interdiscursividade que permitam a identificação e compreensão dos diferentes posicionamentos e/ou perspectivas em jogo, do papel da paráfrase e de produções como as \textbf{paródias} e a \textbf{estilizações}.
Reconstrução da textualidade, recuperação e análise da organização textual, da \textbf{progressão} temática e estabelecimento de relações entre as partes do texto {\EmojiFont ➢} Estabelecer relações entre as partes do texto, identificando repetições, substituições e os elementos coesivos que contribuem para a continuidade do texto e sua \textbf{progressão} temática. {\EmojiFont ➢} Estabelecer relações lógico-discursivas variadas ( identificar/distinguir e relacionar fato e opinião ; causa/efeito ; tese/argumentos ; problema/solução ; definição/exemplos etc. ). {\EmojiFont ➢} Selecionar e hierarquizar informações, tendo em vista as condições de produção e recepção dos textos.
Reflexão crítica sobre as temáticas tratadas e validade das informações {\EmojiFont ➢} Refletir criticamente sobre a fidedignidade das informações, as temáticas, os fatos, os acontecimentos, as questões controversas presentes nos textos lidos, posicionando -se.
Compreensão dos efeitos de sentido provocados pelos usos de recursos linguísticos e \textbf{multissemióticos} em textos pertencentes a gêneros diversos {\EmojiFont ➢} Identificar implícitos e os efeitos de sentido decorrentes de determinados usos expressivos da linguagem, da pontuação e de outras notações, da escolha de determinadas palavras ou expressões e identificar efeitos de ironia ou humor. {\EmojiFont ➢} Identificar e analisar efeitos de sentido decorrentes de escolhas e formatação de imagens ( enquadramento, ângulo/vetor, cor, brilho, contraste ), de sua sequenciação ( disposição e transição, movimentos de \textbf{câmera}, remix ) e da performance – movimentos do corpo, gestos, ocupação do espaço cênico e elementos sonoros ( entonação, trilha sonora, sampleamento etc. ) que nela se relacionam. {\EmojiFont ➢} Identificar e analisar efeitos de sentido decorrentes de escolhas de \textbf{volume}, timbre, intensidade, pausas, ritmo, efeitos sonoros, sincronização etc. em artefatos sonoros.
Estratégias e procedimentos de leitura {\EmojiFont ➢} Selecionar procedimentos de leitura adequados a diferentes objetivos e interesses, levando em conta características do gênero e suporte do texto, de forma a poder proceder a uma leitura autônoma em relação a temas familiares. {\EmojiFont ➢} Estabelecer/considerar os objetivos de leitura. {\EmojiFont ➢} Estabelecer relações entre o texto e conhecimentos prévios, vivências, valores e crenças. {\EmojiFont ➢} Estabelecer expectativas ( pressuposições antecipadoras dos sentidos, da forma e da função do texto ), apoiando -se em seus conhecimentos prévios sobre gênero textual, suporte e universo temático, bem como sobre saliências textuais, recursos gráficos, imagens, dados da própria obra ( índice, prefácio etc. ), confirmando antecipações e inferências realizadas antes e durante a leitura de textos. {\EmojiFont ➢} Localizar/recuperar informação. {\EmojiFont ➢} Inferir ou deduzir informações implícitas. {\EmojiFont ➢} Inferir ou deduzir, pelo contexto semântico ou linguístico, o significado de palavras ou expressões desconhecidas. {\EmojiFont ➢} Identificar ou selecionar, em função do contexto de ocorrência, a acepção mais adequada de um vocábulo ou expressão. {\EmojiFont ➢} Apreender os sentidos globais do texto. {\EmojiFont ➢} Reconhecer/inferir o tema. {\EmojiFont ➢} Articular o verbal com outras linguagens – diagramas, ilustrações, fotografias, vídeos, arquivos sonoros etc. – reconhecendo relações de reiteração, complementaridade ou contradição entre o verbal e as outras linguagens. {\EmojiFont ➢} Buscar, selecionar, tratar, analisar e usar informações, tendo em vista diferentes objetivos. {\EmojiFont ➢} Manejar de forma produtiva a não linearidade da leitura de \textbf{hipertextos} e o manuseio de várias janelas, tendo em vista os objetivos de leitura.
Adesão às práticas de leitura {\EmojiFont ➢} Mostrar -se interessado e envolvido pela leitura de \textbf{livros} de literatura, textos de divulgação científica e/ou textos jornalísticos que circulam em várias mídias. {\EmojiFont ➢} Mostrar -se ou tornar -se receptivo a textos que rompam com seu universo de expectativa, que representem um desafio em relação às suas possibilidades atuais e suas experiências anteriores de leitura, apoiando -se nas marcas linguísticas, em seu conhecimento sobre os gêneros e a temática e nas orientações dadas pelo professor.
FONTE : BNCC, 2017.

% matched lemmas: comentar, câmera, estilização, hipertexto, livro, multiletramento, multissemiótico, paródia, progressão, volume
\end{textsample}
